\section{Higher-form symmetries of Maxwell theory}
\label{sec:maxwell}

Our treatment of linear gravity will follow and generalise that for spin-1, so we begin with a brief review of Maxwell theory on a $d$-dimensional manifold $\mathcal{M}$ with Lorentzian signature. The curvature 2-form $F$ can be written locally in terms of a $U(1)$ connection 1-form $A$, with $F=\dd{A}$. The action is
\begin{equation}\label{eq:S_Maxwell}
    S_{\text{M}} = -\frac{1}{2e^2} \int_{\mathcal{M}} F \wedge \star F
\end{equation}
where $e$ is a coupling constant.
The Bianchi identity $\dd{F} = 0$ and the field equation $\dd\star F=0$ imply the existence of two conserved 2-form currents
\begin{equation}
    J_{\text{e}} = \frac{1}{e^2} F\qc J_{\text{m}} = \frac{1}{2\pi} \star F
\end{equation}
These currents can  be integrated to give
 topological operators 
\begin{equation}\label{eq:Q_m}
    Q_{\text{e}} = \int_{\Sigma_{d-2}} \star J_e\qc Q_{\text{m}} = \int_{\Sigma_2} \star J_m
\end{equation}
for any codimension-2 cycle $\Sigma_{d-2}$ and dimension-2 cycle $\Sigma_2$.
The charge $Q_{\text{e}}$ generates a $U(1)$ symmetry acting on Wilson lines and the charge  $Q_{\text{m}}$ generates a $U(1)$ symmetry acting on 't Hooft lines in $d=4$, or 't Hooft $(d-3)$-branes in $d>4$ \cite{Gaiotto2015GeneralizedSymmetries}.

The field strength $F$ is invariant under shifts of the gauge field $A\to A+\lambda$ by a 1-form $\lambda$ which is closed, $\dd\lambda=0$. 
This is then a symmetry of the action associated with the conserved current $J_{\text{e}}$.
The global symmetries of this type are those which are not gauge symmetries, which means that $\lambda $ is closed but not exact, so that the symmetry group is given by the first cohomology group.
In fact, the symmetry is slightly larger: the gauge field can be shifted by a flat connection modulo gauge transformations \cite{Gaiotto2015GeneralizedSymmetries}.

The theory can also be formulated in terms of a dual $(d-3)$-form potential $\tilde{A}$ with $F=\star \dd \tilde{A}$. There is then a $(d-3)$-form shift symmetry $\tilde{A}\to \tilde{A} + \tilde{\lambda}$ with $\dd\tilde \lambda=0$ associated with the conserved current $J_{\text{m}}$. 
Again, the non-trivial transformations correspond to closed forms $\tilde{\lambda}$ modulo exact ones.
As before, there is a larger symmetry in which $\tilde{A}$ is shifted by a flat gerbe connection modulo gauge transformations.

Note that, in the standard formulation of the theory in terms of the 1-form gauge field $A$, the `magnetic' $(d-3)$-form symmetry does not correspond to any transformation of the gauge field $A$ and in the dual description in terms of the $(d-3)$-form gauge field $\tilde{A}$, the `electric' 1-form symmetry does not correspond to a transformation of the gauge field $\tilde{A}$.

These two higher-form symmetries are viewed as global symmetries, as the parameters are restricted to be closed, and have a mixed 't Hooft anomaly \cite{Gaiotto2015GeneralizedSymmetries}. This can be seen by coupling  to background gauge fields. The electric 1-form symmetry can be gauged by introducing a 2-form background $U(1)$ gauge field $B$ which transforms as $B \to B + \dd\lambda$. This implies that the combination $F-B$ is invariant, and the gauging is achieved by replacing $F$ in \eqref{eq:S_Maxwell} with $F-B$,
\begin{equation}\label{eq:Maxwell_electric_gauged}
    S_{\text{M}}^{\text{e}} = -\frac{1}{2e^2} \int_{\mathcal{M}} (F-B)\wedge\star(F-B)
\end{equation}
 Since the magnetic symmetry is not related to any transformation of the gauge field $A$, it can be gauged by introducing the standard coupling of the current $J_{\text{m}}$ to a background $(d-2)$-form $U(1)$ gauge field $\tilde{B}$,
\begin{equation}
    S_{\text{M}}^{\text{m}} = -\frac{1}{2e^2} \int_{\mathcal{M}} F\wedge\star F + \frac{1}{2\pi} \int_{\mathcal{M}} F \wedge \tilde{B}
\end{equation}
where $\tilde{B} \to \tilde{B} + \dd{\tilde{\lambda}}$ under the magnetic symmetry.

Introducing both the electric and magnetic background fields results in
\begin{equation}
    S_{\text{M}}^{\text{e}+\text{m}} = -\frac{1}{2e^2} \int_{\mathcal{M}} (F-B)\wedge\star(F-B) + \frac{1}{2\pi} \int_{\mathcal{M}} F \wedge \tilde{B}
\end{equation}
This is invariant under the magnetic transformation $\tilde{B} \to \tilde{B} + \dd{\tilde{\lambda}}$, but is no longer invariant under the electric transformations $A\to A+\lambda$, $B\to B+\dd{\lambda}$, signalling an anomaly. In particular, we find
\begin{equation}\label{eq:Maxwell_anomaly}
    \delta S_{\text{M}}^{\text{e}+\text{m}} = \frac{1}{2\pi} \int_{\mathcal{M}} \dd{\lambda} \wedge \tilde{B}
\end{equation}
A counterterm $-\frac{1}{2\pi} B\wedge\tilde{B}$ can be added to the Lagrangian to cancel this variation, but the resulting action is then not invariant under the magnetic transformation. 
There is no counterterm that gives an action that is invariant under both symmetries. 

The anomaly can be cancelled by the inflow mechanism if it is coupled to a $(d+1)$-dimensional theory. This involves extending $B$ and $\tilde{B}$ into a $(d+1)$-dimensional bulk $\mathcal{N}$, with $\partial\mathcal{N} = \mathcal{M}$, and coupling to a topological field theory with a BF action
\begin{equation}\label{eq:S_inflow_Maxwell}
    S_{\text{M}}^{\text{inflow}} = -\frac{1}{2\pi} \int_{\mathcal{N}} B \wedge \dd\tilde{B}
\end{equation}
This $(d+1)$-dimensional theory is non-trivial, indicating that there is no choice of $d$-dimensional counterterm which will remove the anomaly. 
Alternatively, we can understand the anomaly via the descent procedure. The $(d+2)$-form anomaly polynomial which produces the $(d+1)$-dimensional  action \eqref{eq:S_inflow_Maxwell} by descent is 
\begin{equation}\label{eq:anomaly_poly_maxwell}
    \mathcal{I}_{d+2}^{\text{M}} = -\frac{1}{2\pi} \dd{B} \wedge \dd{\tilde{B}}
\end{equation}

A mixed 't Hooft anomaly between two symmetries implies that gauging one explicitly breaks the other. In terms of the topological operators $Q_\text{e}$ and $Q_\text{m}$, this means that gauging the electric symmetry, for example, will render the magnetic operator $Q_\text{m}$ non-topological. 
While the Bianchi identity $\dd{F}=0$ remains valid in the gauged theory, $F$ is no longer a gauge-invariant operator. The gauge-invariant quantity of interest is $F-B$, which is not closed for generic background field configurations since $\dd{(F-B)} = -\dd{B} \neq0$. Therefore, $Q_m$ as defined in \eqref{eq:Q_m} is not a gauge-invariant operator in the gauged theory, and modifying its definition to $Q_m = \frac{1}{2\pi} \int_{\Sigma_2} (F-B)$ results in an operator which is gauge-invariant but not topological. The magnetic symmetry is, therefore, broken when the electric symmetry is coupled to a background gauge field. 

For further discussion of the generalised symmetries and anomaly in Maxwell theory, see e.g.\ the reviews in \cite{Hofman:2017vwr, Brennan:2023mmt, Schafer-Nameki:2023jdn, Bhardwaj:2023kri} and references therein.